% Estilo compartido. Compilar el volumen principal con XeLaTeX o Tectonic.
\documentclass[12pt,letterpaper]{article}
\usepackage[spanish,es-noshorthands]{babel}
\usepackage[left=23mm,right=23mm,top=23mm,bottom=22mm,headheight=15pt]{geometry}
\usepackage{graphicx,xcolor,fancyhdr}
\usepackage[unicode,pdflang={es-MX},bookmarksopen,colorlinks=true,linkcolor=azul]{hyperref}
% La familia Latin Modern de 10 puntos también cubre los títulos grandes.
\DeclareFontShape{TU}{lmr}{m}{n}{<->\UnicodeFontFile{lmroman10-regular}{\UnicodeFontTeXLigatures}}{}
\DeclareFontShape{TU}{lmr}{b}{n}{<->\UnicodeFontFile{lmroman10-bold}{\UnicodeFontTeXLigatures}}{}
\DeclareFontShape{TU}{lmr}{bx}{n}{<->\UnicodeFontFile{lmroman10-bold}{\UnicodeFontTeXLigatures}}{}
\DeclareFontFamily{TU}{lmtt}{}
\DeclareFontShape{TU}{lmtt}{m}{n}{<->\UnicodeFontFile{lmmono10-regular}{}}{}
\definecolor{azul}{HTML}{075479}
\definecolor{tinta}{HTML}{172F3D}
\definecolor{muted}{HTML}{455E6A}
\definecolor{niebla}{HTML}{EDF4F6}
\renewcommand{\headrulewidth}{0pt}
\setlength{\parindent}{0pt}
\setlength{\parskip}{4pt}
\setlength{\emergencystretch}{1.3em}
\linespread{1.08}
\pagestyle{fancy}\fancyhf{}
\fancyhead[L]{\small\color{muted}ESCOM \enspace / \enspace Trabajo Terminal 2026-B155}
\fancyhead[R]{\small\color{muted}\volumen}
\fancyfoot[L]{\small\hyperlink{indice}{Volver al índice}}
\fancyfoot[R]{\small\color{muted}\thepage}
\newcommand{\boton}[1]{\textbf{#1}}
\newcommand{\tema}[1]{\par\addvspace{9pt}{\large\bfseries\color{azul}#1\par}\nobreak}
\newcommand{\pagina}[3]{\clearpage\phantomsection\addcontentsline{toc}{section}{#2}%
{\small\bfseries\color{azul}\MakeUppercase{#1}\par}\vspace{2pt}%
{\raggedright\fontsize{23}{26}\selectfont\bfseries\color{tinta}#2\par}\nobreak\vspace{2pt}%
{\large\color{muted}#3\par}\vspace{4pt}}
\newenvironment{pasos}{\begin{enumerate}\setlength{\itemsep}{3pt}\setlength{\topsep}{1pt}\setlength{\parskip}{0pt}}{\end{enumerate}}
\newcommand{\nota}[2]{\par\addvspace{7pt}\noindent\begingroup\setlength{\fboxsep}{8pt}%
\colorbox{niebla}{\parbox{\dimexpr\linewidth-2\fboxsep\relax}{\small\color{tinta}\textbf{#1}\enspace #2}}\endgroup\par}
\newcommand{\captura}[4][1]{\par\addvspace{6pt}\begin{center}\setlength{\fboxsep}{0pt}%
\fcolorbox{niebla}{white}{\includegraphics[width=\dimexpr#1\linewidth-2\fboxrule\relax,height=#3,keepaspectratio]{figuras/#2.png}}%
\par\vspace{2pt}{\small\color{muted}#4\par}\end{center}}
\newcommand{\portada}[2]{%
\hypersetup{pdftitle={Manual de usuario. #1},pdfauthor={Díaz Ramírez Daniel y Téllez De La Cruz Esaul},pdfsubject={Aplicación web para la gestión de actividades de la ESCOM. TT 2026-B155}}
\color{tinta}\thispagestyle{empty}
\begin{minipage}[c]{.18\linewidth}\includegraphics[width=.83\linewidth]{figuras/escudo-escom.png}\end{minipage}\hfill
\begin{minipage}[c]{.79\linewidth}{\large\bfseries Instituto Politécnico Nacional\par}Escuela Superior de Cómputo\par Trabajo Terminal 2026-B155\end{minipage}
\par\vspace{14mm}{\color{azul}\bfseries MANUAL DE USUARIO\par}\vspace{1mm}
{\fontsize{32}{35}\selectfont\bfseries #1\par}\vspace{2mm}
{\large Aplicación web para la gestión de actividades de la ESCOM\par}\vspace{4mm}
#2\par\vspace{4mm}
\hypertarget{indice}{}\tableofcontents
\vfill
{\small\color{muted}\textbf{Autores}\quad Díaz Ramírez Daniel y Téllez De La Cruz Esaul\par
\textbf{Directora}\quad Dra. Mónica Rivera De La Rosa\par
Edición de noviembre de 2026 \enspace · \enspace Ciudad de México\par}}
\AtBeginDocument{\renewcommand{\contentsname}{Índice de tareas}}

\newcommand{\volumen}{Operación institucional}
\begin{document}
\portada{Operación institucional}{Para el alumno responsable de club, el profesor coordinador, las oficinas de la ESCOM y las cuentas de administración. Entra con la cuenta asignada a tu función y elige la tarea que te corresponde.}

\pagina{Comunicación Social y administración}{Publicar un comunicado}{Abre Avisos con la cuenta de Comunicación Social o de administración. El área emisora identifica la procedencia del mensaje y puede ser distinta de la oficina que lo publica.}
\captura{comunicado-form}{51mm}{La ventana de publicación permite redactar el comunicado o preparar un borrador desde un mensaje.}
\tema{Preparar y publicar}
\begin{pasos}
\item En \boton{Avisos}, selecciona \boton{Publicar comunicado}. Escribe el titular y el cuerpo. Si partes de un oficio, despliega \boton{Asistente IA desde mensaje u oficio}, pega el texto y pulsa \boton{Estructurar con IA}. Revisa nombres, fechas, requisitos y enlaces en el borrador. Si falla el asistente, completa los campos a mano.
\item Selecciona categoría y área emisora. Revisa la \boton{Fecha de publicación}. Para programarlo, elige la fecha y hora futuras. Marca \boton{Visible para visitantes sin sesión} si el aviso es público. Puedes destacarlo, agregar un enlace oficial o adjuntar el cartel.
\item Revisa las versiones para redes. Puedes editar y copiar el texto para publicarlo desde la cuenta de cada red. Pulsa \boton{Publicar comunicado} para guardar el aviso en la plataforma.
\end{pasos}
\tema{Verificar y corregir}
Localiza el titular en \boton{Avisos}. Si es público e inmediato, comprueba también su visibilidad como visitante. Si es futuro, consulta los programados desde la vista de gestión y verifica la fecha antes de crear otro. Para corregirlo, abre la edición, modifica los campos y pulsa \boton{Guardar cambios}. El panel de Comunicación muestra aperturas y clics en los enlaces de los comunicados.

\pagina{Ventanillas y moderación de área}{Responder dudas y mantener las guías}{Gestión Escolar, Becas y CATT atienden sus propios trámites desde Foro y Guías. Usa el selector de área y la cuenta que corresponde a la guía o a la pregunta.}
\captura{nueva-guia}{53mm}{Nueva guía interactiva. El área responsable y el contenido del borrador se revisan antes de publicar.}
\tema{Dar una respuesta oficial}
En \boton{Foro y Guías}, filtra por \boton{Abiertas} y por área. Despliega las respuestas del hilo, usa \boton{Aportar respuesta} y publica las indicaciones. Si una respuesta existente resuelve la duda, usa \boton{Marcar como respuesta oficial} cuando esté disponible. Comprueba que aparezca el distintivo \boton{Respuesta oficial} en la publicación.
\tema{Crear o actualizar una guía}
\begin{pasos}
\item Pulsa \boton{Nueva guía}. Escribe el procedimiento o usa \boton{Elegir archivo} para extraer un borrador desde PDF, DOCX, texto o imagen, hasta 15 MB. Si la extracción es incompleta, corrige el contenido en el editor.
\item Verifica el área, el título y el resumen. Añade los pasos con sus instrucciones. En cada decisión indica la opción y su paso de destino. Recorre las alternativas, revisa las preguntas frecuentes y adjunta el formato PDF público si corresponde. Publica la guía y búscala por su título.
\item Para actualizarla, abre \boton{Abrir guía y mini-foro} y elige \boton{Editar guía visualmente}. Revisa el historial de versiones tras guardar. Para una sugerencia pendiente, lee el paso afectado, ajusta \boton{Texto actualizado del paso} y pulsa \boton{Aceptar e integrar} o \boton{Rechazar}. Si otra persona modificó la guía, vuelve a leer la versión vigente antes de resolver.
\end{pasos}
\par En la ficha completa, consulta las analíticas de uso y cambia el período para revisar la actividad de la guía.

\pagina{Responsables estudiantiles y docentes}{Gestionar un club activo}{El alumno presidente y el profesor coordinador operan su club desde la misma ficha que consultan los integrantes. Los controles adicionales aparecen al entrar con la cuenta designada.}
\captura{club-gestion}{45mm}{Cabecera de un club visto por su responsable. Las pestañas dan acceso a solicitudes, integrantes y métricas.}
\tema{Revisar solicitudes de ingreso}
En \boton{Clubes}, abre \boton{Ver club} y la pestaña \boton{Solicitudes}. Revisa los datos del aspirante y sus respuestas a las preguntas de admisión. Pulsa \boton{Aceptar} o \boton{Rechazar} en esa solicitud. Si la aceptaste, comprueba que la persona aparezca en \boton{Integrantes}. El área autorizada también puede revisar solicitudes dentro de su ámbito.
\tema{Publicar para los integrantes}
Como presidente o coordinador, pulsa \boton{Nueva publicación}, escribe el título y el mensaje y añade una imagen si la necesitas. Publica y comprueba el contenido en \boton{Publicaciones}. Para actualizar horario, contactos o preguntas de ingreso, abre \boton{Configuración}, edita la ficha y guarda. Verifica el resultado en \boton{Acerca y contacto}.
\tema{Proponer una actividad}
En la ficha del club, elige \boton{Nuevo evento} o \boton{Proponer evento}. Captura el título, la fecha, la sede y los detalles del encuentro. Envía la propuesta y revisa su estado en \boton{Eventos}. Una propuesta pendiente debe recibir resolución del área correspondiente. Una vez publicada, consulta en \boton{Métricas} el interés y las personas que marcaron \boton{Asistiré}.
\nota{Si falta un control}{Comprueba el nombre del club y tu cuenta en la cabecera. La membresía aceptada permite ver las publicaciones internas. La presidencia o la coordinación asignadas habilitan la publicación y la propuesta de eventos.}

\pagina{Áreas responsables de clubes}{Revisar y registrar un club}{Subdirección Académica y Clubes Académicos atienden iniciativas académicas. La SSEIS, el área de Clubes Culturales y Deportivos y las personas con designación vigente atienden los tipos que les corresponden.}
\tema{Resolver una propuesta enviada por estudiantes}
\begin{pasos}
\item En \boton{Clubes}, abre \boton{Panel del área} y elige una de las \boton{Solicitudes formales de nuevos clubes}. Verifica el tipo de club, el alumno presidente, el coordinador confirmado y el proyecto de actividades.
\item Abre los expedientes individuales. Revisa carrera, contacto, credencial, boleta o tira de materias, firma y vigencia de derechos IMSS. Registra la observación sobre el archivo que deba corregirse. Conserva la revisión física que corresponda al procedimiento de tu oficina.
\item En la junta conjunta propone opciones de horario, modalidad y lugar o enlace. Consulta las disponibilidades del presidente y del coordinador. Confirma la fecha definitiva y, una vez celebrada, registra que la junta se completó.
\item En \boton{Alta institucional del club}, pulsa \boton{Aprobar y registrar club}. Abre \boton{Ver el club registrado} y comprueba la ficha. Si procede devolver o rechazar la solicitud, usa el dictamen y explica qué falta o cuál es el motivo.
\end{pasos}
\tema{Dar de alta un club ya reconocido}
Desde \boton{Panel del área}, pulsa \boton{+ Registrar club}. Prepara dos PDF de hasta 10 MB cada uno. Adjunta la solicitud escrita con proyecto ejecutivo y el padrón documental de al menos 20 estudiantes. Captura objetivo, cronograma anual, actividades, espacios y contribución. Selecciona una cuenta activa de alumno presidente regular de tercero a séptimo semestre y una cuenta activa de profesor coordinador. Completa los datos, registra el club y búscalo en \boton{Clubes activos}.
\tema{Designar seguimiento en la SSEIS}
La titularidad de la SSEIS o administración puede usar \boton{Designar responsable} desde el panel. Selecciona a la persona, agrega una nota o cargo si procede y pulsa \boton{Designar responsable}. Verifica la designación en el listado. Cuando deje de corresponder, selecciona \boton{Retirar designación} en su registro.

\pagina{Sedes y Prefectura}{Publicar eventos y mantener el mapa}{Becas atiende el Auditorio y Subdirección Académica la Sala 14. Administración tiene acceso a las sedes. Prefectura y administración pueden editar las fichas base del mapa.}
\captura{eventos}{51mm}{Formulario de evento autorizado. La sede determina qué oficina puede publicarlo y si exige ponente.}
\tema{Incorporar una actividad a la cartelera}
\begin{pasos}
\item Una vez autorizada la actividad por el plantel, entra en \boton{Eventos} y pulsa \boton{+ Publicar evento}. Selecciona el espacio. Escribe el título y el nombre del ponente, obligatorio en Auditorio y Sala 14.
\item Registra inicio y, si corresponde, fin posterior al inicio. Añade descripción y cartel. Activa \boton{Abierto al público} si la actividad debe verse en la cartelera pública. Pulsa \boton{Publicar evento autorizado}.
\item Busca la fecha y filtra por la sede para comprobar el registro. Si aparece un conflicto, revisa el evento señalado y confirma con la oficina el horario correcto antes de reenviar.
\end{pasos}
Para propuestas de clubes en Auditorio o Sala 14, la oficina de la sede revisa los pendientes y decide si aprueba y publica o rechaza. Para otras sedes, la revisión corresponde al área del tipo de club. La reservación del espacio se tramita por el canal institucional.
\tema{Corregir un espacio del directorio}
En \boton{Mapa}, busca el código y selecciona el espacio. Abre \boton{Editar ficha base del espacio}, verifica tipo, nombre o función, detalles y profesores confirmados. Pulsa \boton{Guardar ficha}. Vuelve a seleccionar el espacio y revisa que la ficha y el directorio reflejen la corrección.

\pagina{Convivencia y administración}{Registrar infracciones y resolver apelaciones}{Accede a Moderación desde la cabecera con una cuenta habilitada. Las oficinas con navegación restringida atienden sus trámites desde sus propios módulos.}
\captura{moderacion}{42mm}{El historial distingue infracciones, suspensiones y apelaciones. La escala está visible en el panel.}
\tema{Registrar una infracción}
\begin{pasos}
\item Pulsa \boton{Registrar infracción}. Localiza a la persona por nombre, correo o boleta según las opciones disponibles y confirma la identidad antes de seleccionarla.
\item Revisa el área, elige el origen del incidente y describe el hecho concreto y el motivo. Registra la infracción y verifica el nuevo asiento en \boton{Infracciones}.
\item Si se alcanzó un umbral, consulta \boton{Suspensiones}. La escala acumulada es de un mes a las tres infracciones, tres meses a las seis, seis meses a las nueve y suspensión permanente a partir de doce.
\end{pasos}
\tema{Resolver una apelación}
Con cuenta de administración, abre \boton{Apelaciones} en el módulo de moderación. Lee la explicación de la persona y el antecedente de suspensión. Escribe la resolución y su motivo, y acepta o rechaza la apelación. Comprueba el estado y la nota registrada. Esa explicación se muestra a la persona en su consulta de apelaciones.
\tema{Cerrar un reporte falso o duplicado}
Con cuenta de administración, abre la ficha en \boton{Objetos}, localiza \boton{Cerrar reporte falso o spam}, registra el motivo y confirma el cierre. Comprueba el estado \boton{Cerrado}. La entrega de un objeto a su dueño se registra entre las partes mediante el QR descrito en el manual de comunidad.

\pagina{Administración y superadministración}{Gestionar cuentas y permisos}{Con una cuenta de administración, pulsa Panel en la cabecera. Las pestañas del Panel de operación reúnen usuarios, padrón, solicitudes de privacidad, auditoría e indicadores.}
\captura{panel}{39mm}{La pestaña Usuarios y roles permite buscar una cuenta y abrir sus controles.}
\tema{Localizar una cuenta}
En \boton{Usuarios y roles}, busca por nombre, boleta, correo o área. Usa los filtros de rol y estado si hay coincidencias. En la fila correcta pulsa \boton{Gestionar}. Antes de guardar cualquier cambio, contrasta nombre y correo con los de la persona solicitada.
\tema{Asignar un rol institucional}
Esta tarea requiere superadministración. En la ficha de gestión selecciona \boton{Rol}. Para moderación también selecciona el \boton{Área}, por ejemplo Becas, CATT o Comunicación Social. Pulsa \boton{Guardar rol}. Revisa el rol y el área en el listado. Pide a la persona que vuelva a entrar para comprobar su navegación.
\tema{Aplicar o retirar un veto}
Sobre una cuenta que admita la acción, escribe un motivo de al menos diez caracteres y pulsa \boton{Vetar y cerrar sesión}. Para retirarlo, busca la cuenta con el filtro \boton{Vetados}, abre su ficha, justifica el cambio y pulsa \boton{Quitar veto}. Comprueba el estado y consulta el registro en \boton{Auditoría}.
\tema{Ajustar parámetros}
En \boton{Configuración}, revisa el umbral de votos de las propuestas y el intervalo en minutos entre preguntas del foro. Captura el valor acordado y pulsa \boton{Guardar parámetros}. Vuelve a abrir la pestaña para verificarlo. Los requisitos documentales del registro institucional de clubes se revisan en su expediente.

\pagina{Administración}{Invitar cuentas mediante un padrón}{Usa Padrón CSV para enviar invitaciones a estudiantes y profesores que aún no tienen cuenta. Necesitas una lista de nombres y correos institucionales.}
\captura{padron}{57mm}{Padrón CSV ofrece una plantilla vacía, un ejemplo ficticio y un selector de archivo.}
\tema{Preparar el archivo}
Descarga la plantilla desde \boton{Panel} y \boton{Padrón CSV}. Conserva los encabezados \texttt{nombre,correo} y escribe una persona por fila. Exporta como CSV separado por comas. El límite es de 200 filas de personas y 2 MB por archivo. El dominio determina el rol inicial de alumno o profesor.
\tema{Previsualizar, confirmar y revisar}
\begin{pasos}
\item Selecciona el CSV en \boton{Archivo CSV}. Revisa cuántas filas están listas y cuántas presentan error. Seleccionar el archivo abre la vista previa.
\item Si hay errores, usa \boton{Descargar errores CSV}, corrige los nombres o correos señalados, elimina duplicados y vuelve a seleccionar el archivo corregido. La importación se habilita al resolver los errores de la vista previa.
\item Verifica el total y pulsa \boton{Confirmar importación}. Revisa el resultado por fila. Una invitación creada requiere que su destinatario abra el mensaje, confirme la invitación y elija su contraseña.
\item Si hay errores después del envío, descarga las filas afectadas. Revisa las cuentas ya existentes y prepara un nuevo archivo únicamente con los casos pendientes que sí necesitan invitación.
\end{pasos}
\nota{Privilegios}{La asignación de moderadores, administradores y áreas se realiza desde \textbf{Usuarios y roles} con superadministración. La plantilla de invitación solo requiere nombre y correo.}

\pagina{Administración}{Atender solicitudes y obtener reportes}{El panel permite responder a las personas sobre sus datos, consultar quién realizó una acción y exportar indicadores para un informe.}
\tema{Dar seguimiento a una solicitud de privacidad}
\begin{pasos}
\item En \boton{Panel}, abre \boton{Derechos ARCO}. Localiza la solicitud pendiente y lee el tipo y la descripción. ARCO reúne los derechos de acceso, rectificación, cancelación y oposición.
\item Escribe una nota para el titular. Usa \boton{En revisión} para registrar que la estás atendiendo. Tras realizar la gestión correspondiente, escribe la resolución y pulsa \boton{Resolver} o \boton{Rechazar}.
\item Verifica el nuevo estado. La persona puede consultar la respuesta en su perfil, en \boton{Mis solicitudes registradas}.
\end{pasos}
\tema{Localizar una acción y exportar su registro}
En \boton{Auditoría}, captura el nombre del actor en \boton{Usuario} o usa una entidad, por ejemplo \texttt{clubs} o \texttt{avisos}. Acota \boton{Desde} y \boton{Hasta} y pulsa \boton{Filtrar}. Comprueba fecha, actor, entidad y acción. El botón \boton{CSV} descarga el resultado de los filtros aplicados. Si no hay filas, amplía el intervalo o quita un filtro.
\tema{Preparar el informe de impacto}
En \boton{Resumen e impacto}, consulta los indicadores de usuarios, preguntas, clubes, avisos y objetos. Los desgloses de guías y preguntas se muestran al alcanzar cinco lectores o cinco autores, respectivamente. Para descargar los datos agregados de guías y preguntas del foro, usa \boton{Exportar CSV}. Para conservar el informe, pulsa \boton{Imprimir o guardar PDF}, elige guardar como PDF en el diálogo del navegador y revisa la vista previa.
\nota{Qué archivo elegir}{Para analizar lecturas de guías y preguntas por área, usa el CSV de impacto. Para documentar quién realizó una acción concreta y cuándo, exporta el CSV filtrado de auditoría.}
\end{document}
